% ============================================================
% PowerSystemBook 3.2 -- central preamble
% Engine: LuaLaTeX | Bibliography: Biber | Index: MakeIndex
% ============================================================

% 1. Language, fonts, typography
\usepackage[english]{babel}
\usepackage{csquotes}
\usepackage{fontspec}
\usepackage{unicode-math}
\usepackage{microtype}
\setmainfont{Libertinus Serif}
\setsansfont{Libertinus Sans}
\setmonofont{Libertinus Mono}[Scale=MatchLowercase]
\setmathfont{Libertinus Math}
\setlength{\emergencystretch}{2em}
\raggedbottom

% 2. Colour palette
\usepackage{xcolor}
\definecolor{BookBlue}{RGB}{24,72,120}
\definecolor{BookBlueDark}{RGB}{17,50,84}
\definecolor{BookGreen}{RGB}{42,105,70}
\definecolor{BookOrange}{RGB}{181,105,22}
\definecolor{BookRed}{RGB}{155,45,45}
\definecolor{BookGray}{RGB}{88,92,96}
\definecolor{LightBlue}{RGB}{241,247,252}
\definecolor{LightGreen}{RGB}{242,249,244}
\definecolor{LightYellow}{RGB}{255,249,230}
\definecolor{LightRed}{RGB}{253,242,242}
\definecolor{LightGray}{RGB}{246,247,248}
% Formula-sheet colours (formerly defined locally in the appendices).
\definecolor{navy}{HTML}{17365D}
\definecolor{blue}{HTML}{2F75B5}
\definecolor{sky}{HTML}{EAF3FA}
\definecolor{gold}{HTML}{D69E2E}
\definecolor{ink}{HTML}{243447}
\definecolor{muted}{HTML}{617285}
\definecolor{shadecolor}{HTML}{F1F6FA}

% 3. Mathematics and units
\usepackage{amsmath,mathtools,bm}
\usepackage{siunitx}
\DeclareSIUnit{\varunit}{var}
\DeclareSIUnit{\rpm}{rpm}
\DeclareSIUnit{\pu}{p.u.}
\sisetup{
  per-mode=symbol,
  separate-uncertainty=true,
  detect-all,
  exponent-product=\cdot,
  output-decimal-marker={.},
  group-minimum-digits=4
}
\newcommand{\vect}[1]{\bm{#1}}
\newcommand{\mat}[1]{\bm{#1}}
\newcommand{\e}{\mathrm{e}}
\newcommand{\imag}{\mathrm{j}}
\newcommand{\diff}{\mathop{}\!\mathrm{d}}
\newcommand{\Rmag}{\mathcal{R}}
\newcommand{\sign}{\operatorname{sign}}
\newcommand{\eqsep}{\addlinespace[2pt]}
\newcommand{\formulaTableStretch}{1.18}

% 4. Graphics, plots, captions, tables
\usepackage{graphicx}
\usepackage{tikz}
\graphicspath{{figures/}{figures/common/}{figures/transformers/}{figures/lines/}{figures/generators/}}
\usepackage{float}
\usepackage{subcaption}
\usepackage{booktabs,longtable,tabularx,array,multirow}
\usepackage{placeins}
\usepackage{pdflscape}
\usepackage{framed}

% Full-page landscape figure.
% Usage:
% \begin{widefigure}{Caption text}{fig:label}
%   \includegraphics[width=.95\linewidth]{drawing.pdf}
% \end{widefigure}
\NewDocumentEnvironment{widefigure}{O{p} m m}{%
  \begin{landscape}%
  \begin{figure}[#1]%
  \centering
}{%
  \caption{#2}%
  \label{#3}%
  \end{figure}%
  \end{landscape}%
}
\captionsetup{
  font=small,
  labelfont={bf,color=BookBlueDark},
  textfont=small,
  justification=raggedright,
  singlelinecheck=false,
  skip=6pt
}
\renewcommand{\arraystretch}{1.15}
\newcolumntype{Y}{>{\raggedright\arraybackslash}X}

% 5. Lists
\usepackage{enumitem}
\setlist{topsep=4pt,itemsep=2pt,parsep=0pt}
\setlist[itemize,1]{label=\textcolor{BookBlue}{\small\textbullet}}

% 6. Page design, headings, headers and footers
\usepackage[automark,headsepline=.4pt]{scrlayer-scrpage}
\setkomafont{pageheadfoot}{\small\sffamily\color{BookGray}}
\setkomafont{pagenumber}{\bfseries\color{BookBlueDark}}
\clearpairofpagestyles
\ihead{\headmark}
\ohead{\pagemark}
\ifoot{Friedrich Schmidt}
\ofoot{Power System Simulation with LTspice}
\pagestyle{scrheadings}
\renewcommand*{\chapterpagestyle}{plain}

\setkomafont{chapter}{\normalfont\sffamily\bfseries\color{BookBlueDark}}
\setkomafont{section}{\normalfont\sffamily\bfseries\color{BookBlue}}
\setkomafont{subsection}{\normalfont\sffamily\bfseries\color{BookBlueDark}}
\setkomafont{disposition}{\normalfont\sffamily\bfseries}
\RedeclareSectionCommand[beforeskip=0pt,afterskip=1.4\baselineskip]{chapter}
\RedeclareSectionCommand[beforeskip=1.7\baselineskip,afterskip=.55\baselineskip]{section}
\RedeclareSectionCommand[beforeskip=1.3\baselineskip,afterskip=.4\baselineskip]{subsection}

% 7. Hyperlinks and cross-references
\usepackage{hyperref}
\hypersetup{
  colorlinks=true,
  linkcolor=BookBlueDark,
  citecolor=BookGreen,
  urlcolor=BookBlue,
  pdfborder={0 0 0},
  pdftitle={Power System Simulation with LTspice},
  pdfauthor={Friedrich Schmidt},
  pdfsubject={Power-system modelling and simulation}
}
\usepackage[nameinlink,noabbrev,capitalise]{cleveref}

% 8. Instructional boxes
% Robust tcolorbox setup for MiKTeX/LuaLaTeX.
% Deliberately no "most", "skins", "enhanced" or TikZ "fadings" library.
\usepackage{tcolorbox}
\tcbuselibrary{breakable,theorems}

% Common appearance: rounded corners, integrated coloured title bar,
% thin frame, generous spacing and automatic page breaks.
\tcbset{
  bookbox/.style={
    breakable,
    boxrule=.55pt,
    arc=1.6mm,
    outer arc=1.6mm,
    left=3.5mm,
    right=3.5mm,
    top=3.0mm,
    bottom=3.0mm,
    before skip=12pt,
    after skip=12pt,
    fonttitle=\sffamily\bfseries,
    coltitle=white,
    separator sign={\enspace--},
    description delimiters parenthesis,
    before upper={\setlength{\parindent}{0pt}},
    titlerule=.4pt
  }
}

\newtcbtheorem[number within=chapter]{definition}{Definition}{
  bookbox,
  colback=LightBlue,
  colframe=BookBlue,
  colbacktitle=BookBlue
}{def}

\newtcbtheorem[number within=chapter]{remark}{Remark}{
  bookbox,
  colback=LightGray,
  colframe=BookGray,
  colbacktitle=BookGray
}{rem}

\newtcbtheorem[number within=chapter]{example}{Example}{
  bookbox,
  colback=LightGreen,
  colframe=BookGreen,
  colbacktitle=BookGreen
}{ex}

\newtcbtheorem[number within=chapter]{important}{Important}{
  bookbox,
  colback=LightRed,
  colframe=BookRed,
  colbacktitle=BookRed
}{imp}

\newtcbtheorem[number within=chapter]{simulation}{Simulation}{
  bookbox,
  colback=white,
  colframe=BookBlue,
  colbacktitle=BookBlue
}{sim}

\newtcbtheorem[number within=chapter]{calculation}{Calculation}{
  bookbox,
  colback=white,
  colframe=BookOrange,
  colbacktitle=BookOrange,
  before upper={
    \setlength{\parindent}{0pt}
    \setlength{\abovedisplayskip}{6pt}
    \setlength{\belowdisplayskip}{6pt}
  }
}{calc}

\newtcbtheorem[number within=chapter]{result}{Result}{
  bookbox,
  colback=LightGreen,
  colframe=BookGreen,
  colbacktitle=BookGreen,
  fontupper=\bfseries
}{res}

\newtcolorbox{chaptersummary}[1][Chapter Summary]{
  bookbox,
  title={#1},
  colback=LightBlue,
  colframe=BookBlueDark,
  colbacktitle=BookBlueDark,
  before upper={
    \setlength{\parindent}{0pt}
    \setlength{\parskip}{.4\baselineskip}
  }
}

\newtcolorbox{warningbox}[1][Warning]{
  bookbox,
  title={#1},
  colback=LightRed,
  colframe=BookRed,
  colbacktitle=BookRed
}

% Formula sheets: lightweight, fading-free design based on the transformer
% appendix.  Both formula appendices use only these central definitions.
\newenvironment{formulaBox}[1]{%
  \par\medskip
  {\sffamily\bfseries\large\color{navy}#1}\par\nobreak\smallskip
  \def\FrameCommand{\hspace{1pt}{\color{blue}\vrule width 2.2pt}\hspace{8pt}}%
  \MakeFramed{\advance\hsize-11.2pt\FrameRestore}%
  \setlength{\parindent}{0pt}\setlength{\parskip}{5pt}%
}{\endMakeFramed\par\smallskip}

\newenvironment{noteBox}{%
  \par\smallskip
  \def\FrameCommand{\hspace{1pt}{\color{gold}\vrule width 2.2pt}%
    \colorbox{shadecolor}}%
  \MakeFramed{\advance\hsize-2\width\advance\hsize-2\fboxsep\FrameRestore}%
  \setlength{\parindent}{0pt}\setlength{\parskip}{5pt}%
}{\endMakeFramed\par\smallskip}

\newcommand{\formulasheettitle}[2]{%
  \refstepcounter{chapter}%
  \addcontentsline{toc}{chapter}{\protect\numberline{\thechapter}#1}%
  \markboth{#1}{Formula sheet}%
  \begingroup
  \setlength{\fboxsep}{8mm}%
  \noindent\colorbox{navy}{%
    \parbox{\dimexpr\textwidth-2\fboxsep\relax}{%
      {\sffamily\bfseries\fontsize{25}{29}\selectfont\color{white}#1}\par
      \vspace{2mm}%
      {\sffamily\fontsize{15}{18}\selectfont\color{white!85}#2}%
    }}%
  \endgroup
  \vspace{3mm}%
}

\crefname{definition}{definition}{definitions}
\crefname{remark}{remark}{remarks}
\crefname{example}{example}{examples}
\crefname{important}{important note}{important notes}
\crefname{simulation}{simulation}{simulations}
\crefname{calculation}{calculation}{calculations}
\crefname{result}{result}{results}

% 9. MATLAB, LTspice and generic code listings
\usepackage{listings}
\renewcommand{\lstlistlistingname}{List of Listings}
\renewcommand{\lstlistingname}{Listing}
\lstdefinestyle{CommonCode}{
  basicstyle=\ttfamily\small,
  numbers=left,numberstyle=\scriptsize\color{BookGray},numbersep=8pt,
  frame=single,framerule=.4pt,rulecolor=\color{BookGray},
  backgroundcolor=\color{LightGray},
  breaklines=true,breakatwhitespace=false,
  showstringspaces=false,tabsize=4,columns=fullflexible,keepspaces=true,
  aboveskip=10pt,belowskip=10pt,
  captionpos=b
}
\lstdefinestyle{Matlab}{style=CommonCode,language=Matlab,keywordstyle=\bfseries\color{BookBlue},commentstyle=\itshape\color{BookGreen},stringstyle=\color{BookOrange}}
\lstdefinestyle{LTspice}{style=CommonCode,numbers=none,numbersep=0pt,morecomment=[l]{*},morecomment=[l]{;},keywordstyle=\bfseries\color{BookBlue}}
\lstdefinestyle{PlainText}{style=CommonCode,numbers=none}

% 10. Bibliography
\usepackage[backend=biber,style=ieee,sorting=none,maxbibnames=99]{biblatex}
\addbibresource{bibliography.bib}

% 11. Index and document depth
\usepackage{imakeidx}
\makeindex[title=Index,intoc]
\setcounter{tocdepth}{2}
\setcounter{secnumdepth}{3}
\listfiles
